%=============================================================================
% DISCUSSION & CONCLUSION
% Rewritten for Phase 3 against results/phase3_final.csv and stage3_search.csv.
%=============================================================================
\section{Discussion and Conclusion}
\label{sec:conclusion}

\subsection{Discussion}
We set out to improve district-level dengue forecasting with a learned graph, a
mechanistic constraint and generative augmentation. One of the three survives
scrutiny, one does not replicate, and one is actively harmful --- and the largest
improvement we measured came from none of them.

\paragraph{Sample size decided more than architecture did.}
The learned adjacency won $17/24$ paired runs at three seeds and $35/64$ at eight.
Nothing about the model changed; only the number of runs did. A paired power
analysis on the three-seed results had already indicated $n \approx 54$ was needed,
which is why the replication was run at all. We think this is the most transferable
lesson here: on a 25-node graph with fold-to-fold RMSE spanning 27 to 68, a search
over configurations at small $n$ will return apparently promising effects at a rate
that has little to do with which of them are real. Two of our four were not.

\paragraph{The missing component was temporal, not spatial.}
Our strongest result is the correction of our own omission. The model contained no
temporal operator --- it flattened the input window and applied a linear map ---
while every method it was compared against carries one. Restoring a gated causal
convolution improved RMSE by 5.0 points, more than any graph or loss modification
we tried, and produced the only model statistically indistinguishable from
persistence. A graph-free LSTM outperforming graph models was the symptom; this was
the cause.

\paragraph{Constraints must be measured before they are weighted.}
The spatial penalty exceeded the data loss by four orders of magnitude, one
constraint term was provably inert --- five runs were bit-identical to the
unconstrained control --- and the usable weight range sat two orders of magnitude
from any value a grid search would naturally cover. The constraint that did work
was found by measuring its gradient magnitude first. We would not have found it by
searching harder.

\paragraph{Naive baselines are not a formality.}
Persistence ($55.12$), a random forest ($56.31$) and a graph-free LSTM ($57.95$)
all outperform the graph models, and the random forest is not significantly worse
than persistence. Prior work on this dataset reports none of these comparisons. We do not conclude
that graph methods cannot help here, but we do conclude that on this benchmark
their advantage has not been demonstrated, and that reporting it requires
comparators that were absent.

\subsection{Limitations}
\textbf{(1) One dataset, one disease.} All findings are for 25 districts of Sri
Lankan dengue; whether they transfer to larger graphs where message passing has
more structure to exploit is untested.
\textbf{(2) Under-ascertainment.} Every model treats reported cases as incidence.
Reporting rates vary by district and season, and no observation model corrects for
it.
\textbf{(3) Point forecasts only.} Calibration is unassessed; CRPS and interval
coverage are not reported.
\textbf{(4) A negative result is not a proof.} We show that these methods do not
help under this protocol, not that no such method could.

\subsection{Future Work}
The mechanistic direction remains open but needs a different formulation. The
SEIR--SEI model has seven compartments of which one is observed, so a hard residual
is under-determined on this data; the constraint that worked here avoided latent
compartments entirely, and extending that family is more promising than
instantiating the full system. Adaptive loss weighting from gradient
statistics~\cite{wang2021gradient} is the natural successor to the fixed weights
and curricula tested here, given how narrow the usable range proved. Beyond that:
probabilistic forecasts evaluated with CRPS and interval
coverage~\cite{bracher2021evaluating}, sparsified graph learning for small
node counts~\cite{wu2020mtgnn}, and validation on
OpenDengue~\cite{clarke2024opendengue} to test whether these conclusions are
specific to this benchmark.
